\documentclass{article}
\usepackage{amsmath}
\usepackage{amssymb}
\usepackage{graphicx}
\usepackage{accessibility}

\title{Low-Dimensional Knot Theory, Artin Braid Groups, and Laurent Skein Invariants}
\author{OxidizeMath Theory Verification}
\date{\today}

\begin{document}

\maketitle

\begin{abstract}
We formalize the mathematical foundations of the \texttt{pure\_math::knot\_theory} and \texttt{algebra} extensions in \texttt{OxidizeMath}: univariate Laurent polynomial rings $R[t, t^{-1}]$, infinite Artin braid groups $B_n$ with smooth 3D Alexander closures, rotation-minimizing Bishop parallel-transport frames for watertight tubular meshes, Gauss writhe and linking integrals, oriented planar link diagrams with Reidemeister I--III ambient isotopy moves, Fox and reduced Burau Alexander polynomials, and state-sum Kauffman and Jones polynomials.
\end{abstract}

\section{Laurent Polynomial Rings}

For a commutative ring $R$, the univariate Laurent polynomial ring $R[t, t^{-1}]$ consists of finite formal sums
\[
P(t) = \sum_{k=m}^{M} c_k t^k, \qquad m, M \in \mathbb{Z}, \; c_k \in R.
\]
Canonical trimming enforces $c_m \neq 0$ and $c_M \neq 0$ for $P \neq 0$, with $P(t)$ called \emph{reciprocal} when $P(t) = P(t^{-1})$. Because the Alexander polynomial is defined up to units $\pm t^k \in \mathbb{Z}[t, t^{-1}]$~\cite{alexander1928topological}, canonical normalization shifts the support to $\lfloor -(M - m)/2 \rfloor$ and chooses the global sign so $P(1) > 0$ (or positive leading coefficient when $P(1) = 0$).

\section{Artin Braid Groups and Alexander Closures}

Emil Artin's braid group $B_n$ ($n \ge 1$) is generated by elementary crossings $\sigma_1, \dots, \sigma_{n-1}$ subject to~\cite{artin1947theory}
\[
\sigma_i \sigma_j = \sigma_j \sigma_i \quad (|i - j| \ge 2), \qquad
\sigma_i \sigma_{i+1} \sigma_i = \sigma_{i+1} \sigma_i \sigma_{i+1} \quad (1 \le i \le n - 2).
\]
By Alexander's theorem, closing the top strand slots concentrically around a solid torus to the bottom slots produces a smooth oriented spatial link $\hat{\beta} \subset \mathbb{R}^3$ whose component count equals the cycle count of the induced strand permutation $\pi(\beta) \in S_n$.

\section{Bishop Parallel-Transport Frames and Gauss Integrals}

Frenet--Serret frames become undefined at inflection points where curvature vanishes. Following Bishop~\cite{bishop1975there}, discrete double-reflection parallel transport propagates an initial unit normal $U_0 \perp T_0$ around a closed curve $\{P_i\}_{i=0}^{N-1}$, distributing the closed-loop holonomy angle $\Delta\psi$ uniformly along arc-index so that the extruded tubular surface
\[
Q_{i,j} = P_i + r\bigl(\cos\phi_j\, U_i + \sin\phi_j\, V_i\bigr), \qquad V_i = T_i \times U_i,
\]
forms a closed, orientable, watertight 2-manifold mesh.

For disjoint oriented closed curves $C_1, C_2 \subset \mathbb{R}^3$, Gauss's linking integral
\[
\operatorname{Lk}(C_1, C_2) = \frac{1}{4\pi} \oint_{C_1} \oint_{C_2}
\frac{(r_1 - r_2) \cdot (dr_1 \times dr_2)}{\|r_1 - r_2\|^3} \in \mathbb{Z}
\]
is a topological invariant evaluated via zero-allocation 2-point Gauss--Legendre quadrature over non-adjacent polyline chords; the analogous principal-value self-integral defines the continuous 3D spatial writhe $W_r(C) \in \mathbb{R}$.

\section{Planar Diagrams and Reidemeister Isotopy}

An oriented planar link diagram $D$ records 4-valent transverse crossings with counter-clockwise incident segment slots $[s_0, s_1, s_2, s_3]$ starting at the incoming under-strand $s_0$, with crossing sign $\varepsilon(c) = \operatorname{sgn}\det(T_{\text{over}}, T_{\text{under}}) \in \{\pm 1\}$ and algebraic writhe $w(D) = \sum_c \varepsilon(c)$. By Reidemeister's theorem, two diagrams represent ambient-isotopic links in $\mathbb{R}^3$ if and only if they are connected by finite sequences of planar isotopies and Reidemeister moves:
\begin{itemize}
    \item \textbf{R-I (twist/untwist):} inserts or removes an adjacent monogon curl, shifting $w(D)$ by $\pm 1$.
    \item \textbf{R-II (poke/unpoke):} inserts or cancels an untwisted opposite-sign bigon, preserving $w(D)$.
    \item \textbf{R-III (triangle slide):} slides an over-strand across an under-crossing, reversing crossing visit order along all three strands while preserving $w(D)$.
\end{itemize}

\section{Alexander, Kauffman, and Jones Polynomials}

From the Wirtinger presentation of $\pi_1(S^3 \setminus L)$, Fox free derivatives assemble a $c \times m$ matrix over $\mathbb{Z}[t, t^{-1}]$ with crossing rows $(1 - t) x_{\text{over}} + t x_{\text{in}} - x_{\text{out}}$ (for $\varepsilon = +1$) and $(1 - t) x_{\text{over}} - x_{\text{in}} + t x_{\text{out}}$ (for $\varepsilon = -1$), whose $(m - 1) \times (m - 1)$ Bareiss minor yields $\Delta_L(t)$~\cite{alexander1928topological}. Equivalently, on $B_n$, the $(n - 1)$-dimensional reduced Burau representation $\bar{\psi} : B_n \to \operatorname{GL}_{n-1}(\mathbb{Z}[t, t^{-1}])$ satisfies
\[
\Delta_{\hat{\beta}}(t) \doteq \frac{\det\bigl(I_{n-1} - \bar{\psi}(\beta)(t)\bigr)}{1 + t + \dots + t^{n-1}}.
\]

Kauffman's state model~\cite{kauffman1987state} expands all $2^c$ smoothings ($c \le 16$) with loop weight $d = -A^2 - A^{-2}$:
\[
\langle D \rangle = \sum_{s} A^{\#A(s) - \#B(s)} d^{|s| - 1}, \qquad
X_L(A) = (-A^3)^{-w(D)} \langle D \rangle,
\]
and the substitution $t^{1/2} = A^{-2}$ ($t = A^{-4}$) yields Vaughan Jones's polynomial $V_L(t)$~\cite{jones1985polynomial}, distinguishing chiral mirror pairs such as the right-handed ($t + t^3 - t^4$) and left-handed ($-t^{-4} + t^{-3} + t^{-1}$) trefoil knots.

\bibliographystyle{plain}
\bibliography{references}

\end{document}
